「弊社はアフリカでの火力発電事業に投資すべきか？」

的確な経営判断を行うためには、対象となる市場や技術について、多角的な観点から情報を収集し、考察を深めることが重要です。
例えば、「アフリカで火力発電のニーズはあるか」「他社はその市場に参入しているか」「現地の規制や政策は事業に追い風となるか」といった観点です。

私は、このようなビジネス特有の観点に基づいて情報を検索し考察するためのアプリケーションを開発しました。
本アプリケーションの中核技術は、テキスト中に現れる関係性を抽出する「関係抽出」技術です。
具体的には、「AはBにニーズがある」「AがBに参入する」「AがBを規制する」「AとBが協業する」といったビジネス上重要な関係性を、構文パターンや述語表現に基づいて解析します。
この技術により、新聞記事やWeb上のデータなど、膨大なテキストデータの中から、経営判断に有用な情報を効率的に検索することができます。

本アプリケーションではさらに、ユーザが入力したテーマに対して、次に調べるべき観点も推薦します。
例えば、ユーザが「X事業に参入すべきか」と考えている場合、「X事業に対するニーズ」「X事業における競合他社の動向」「X事業に関する規制」といった調査観点を提示します。
これにより、単なる情報検索にとどまらず、経営判断に必要な考察を段階的に深めていくことができます。
